% GENERATED FILE. Edit canonical lessons, then run npm run generate:books.
% canonical-source-hash: fnv1a64:38358b35ccc59508
% canonical-lessons: HI-C544-qismat, HI-C544-samyog, HI-C544-mauqa, HI-C544-avsar, HI-C544-dav

\chapter{A Stroke of Luck}
\label{ch:hi-a-stroke-of-luck}
\hlchaptermodality{544}

\begin{chapteropening}
\textbf{By the end of this chapter:} \emph{I can talk about luck and chances.}

Complete the last lesson of chapter 544: I can talk about luck and chances.
\end{chapteropening}

\section[\texorpdfstring{\hi{क़िस्मत} (qismat)}{क़िस्मत (qismat)}]{\hi{क़िस्मत} (qismat) — luck; fortune}

\label{lesson:HI-C544-qismat}



\begin{quote}
Before the new one: say the Hindi for an oath (formal), then the Hindi for bound (by a rule or a time).
\end{quote}

\subsection*{You'll want to know: \hi{क़िस्मत}}
\textbf{\hi{क़िस्मत}} — \emph{qismat} — \textquotedblleft{}luck; fortune\textquotedblright{}.

\begin{quote}
\hi{उसकी} \hi{क़िस्मत} \hi{अच्छी} \hi{है}, \hi{उसे} \hi{हमेशा} \hi{सीट} \hi{मिल} \hi{जाती} \hi{है।}
\end{quote}

\begin{quote}
\emph{Uskī qismat acchī hai, use hameśā sīṭ mil jātī hai.}
\end{quote}

\begin{quote}
He is lucky: he always gets a seat.
\end{quote}

Feminine.

\begin{cousinweb}
\textbf{\hi{क़िस्मत}} (\emph{qismat}) is usually traced to Arabic \emph{qisma}, \textquotedblleft{}a share, a portion\textquotedblright{}, through Persian. One's share, handed out in advance, is one's luck.
\end{cousinweb}

\begin{grammarlens}[title={What you've built}]
One more word for saying what happens if.
\end{grammarlens}

\subsection*{Your turn}
\begin{itemize}
\raggedright
  \item \emph{Say it:} \emph{qismat}
  \item \emph{Say it:} \emph{qismat}, once more
  \item \emph{Say it:} \emph{Uskī qismat acchī hai, use hameśā sīṭ mil jātī hai.} — He is lucky: he always gets a seat.
  \item \emph{Recall:} say the Hindi for an oath (formal), then the Hindi for bound (by a rule or a time), then say \emph{qismat} again
\end{itemize}

\begin{tcolorbox}[breakable,colback=teal!4,colframe=teal!35!black,title={Before you move on}]
What does \hi{क़िस्मत} mean? (\textquotedblleft{}Luck; fortune\textquotedblright{}.) Say it once more.
\end{tcolorbox}
\section[\texorpdfstring{\hi{संयोग} (saṁyog)}{संयोग (saṁyog)}]{\hi{संयोग} (saṁyog) — a coincidence}

\label{lesson:HI-C544-samyog}



\begin{quote}
Before the new one: say the Hindi for bound (by a rule or a time), then the Hindi for luck.
\end{quote}

\subsection*{You'll want to know: \hi{संयोग}}
\textbf{\hi{संयोग}} — \emph{saṁyog} — \textquotedblleft{}a coincidence\textquotedblright{}.

\begin{quote}
\hi{क्या} \hi{संयोग} \hi{है}, \hi{हम} \hi{दोनों} \hi{का} \hi{जन्मदिन} \hi{एक} \hi{ही} \hi{दिन} \hi{है}!
\end{quote}

\begin{quote}
\emph{Kyā saṁyog hai, ham donõ kā janmdin ek hī din hai!}
\end{quote}

\begin{quote}
What a coincidence: we both have our birthday on the same day!
\end{quote}

Masculine.

\begin{cousinweb}
\textbf{\hi{संयोग}} (\emph{saṁyog}) is a tatsama, taken whole from Sanskrit \emph{saṃyoga}, \textquotedblleft{}conjunction, connection\textquotedblright{} (Monier-Williams). Hindi uses it for things that come together by chance.
\end{cousinweb}

\begin{grammarlens}[title={What you've built}]
One more word for saying what happens if.
\end{grammarlens}

\subsection*{Your turn}
\begin{itemize}
\raggedright
  \item \emph{Say it:} \emph{saṁyog}
  \item \emph{Say it:} \emph{saṁyog}, once more
  \item \emph{Say it:} \emph{Kyā saṁyog hai, ham donõ kā janmdin ek hī din hai!} — What a coincidence: we both have our birthday on the same day!
  \item \emph{Recall:} say the Hindi for bound (by a rule or a time), then the Hindi for luck, then say \emph{saṁyog} again
\end{itemize}

\begin{tcolorbox}[breakable,colback=teal!4,colframe=teal!35!black,title={Before you move on}]
What does \hi{संयोग} mean? (\textquotedblleft{}A coincidence\textquotedblright{}.) Say it once more.
\end{tcolorbox}
\section[\texorpdfstring{\hi{मौक़ा} (mauqā)}{मौक़ा (mauqā)}]{\hi{मौक़ा} (mauqā) — a chance; an occasion}

\label{lesson:HI-C544-mauqa}



\begin{quote}
Before the new one: say the Hindi for luck, then the Hindi for a coincidence.
\end{quote}

\subsection*{You'll want to know: \hi{मौक़ा}}
\textbf{\hi{मौक़ा}} — \emph{mauqā} — \textquotedblleft{}a chance; an occasion\textquotedblright{}.

\begin{quote}
\hi{मुझे} \hi{मंच} \hi{पर} \hi{गाने} \hi{का} \hi{मौक़ा} \hi{मिला।}
\end{quote}

\begin{quote}
\emph{Mujhe mañc par gāne kā mauqā milā.}
\end{quote}

\begin{quote}
I got the chance to sing on stage.
\end{quote}

Masculine.

\begin{cousinweb}
\textbf{\hi{मौक़ा}} (\emph{mauqā}) is usually traced to Arabic \emph{mauqiʿ}, \textquotedblleft{}a place where something falls\textquotedblright{}, through Persian. Hindi and Urdu use it for the moment when something can happen.
\end{cousinweb}

\begin{grammarlens}[title={What you've built}]
One more word for saying what happens if.
\end{grammarlens}

\subsection*{Your turn}
\begin{itemize}
\raggedright
  \item \emph{Say it:} \emph{mauqā}
  \item \emph{Say it:} \emph{mauqā}, once more
  \item \emph{Say it:} \emph{Mujhe mañc par gāne kā mauqā milā.} — I got the chance to sing on stage.
  \item \emph{Recall:} say the Hindi for luck, then the Hindi for a coincidence, then say \emph{mauqā} again
\end{itemize}

\begin{tcolorbox}[breakable,colback=teal!4,colframe=teal!35!black,title={Before you move on}]
What does \hi{मौक़ा} mean? (\textquotedblleft{}A chance; an occasion\textquotedblright{}.) Say it once more.
\end{tcolorbox}
\section[\texorpdfstring{\hi{अवसर} (avsar)}{अवसर (avsar)}]{\hi{अवसर} (avsar) — an opportunity (formal)}

\label{lesson:HI-C544-avsar}



\begin{quote}
Before the new one: say the Hindi for a coincidence, then the Hindi for a chance.
\end{quote}

\subsection*{You'll want to know: \hi{अवसर}}
\textbf{\hi{अवसर}} — \emph{avsar} — \textquotedblleft{}an opportunity (formal)\textquotedblright{}.

\begin{quote}
\hi{ऐसा} \hi{अवसर} \hi{रोज़} \hi{नहीं} \hi{आता।}
\end{quote}

\begin{quote}
\emph{Aisā avsar roz nahīṃ ātā.}
\end{quote}

\begin{quote}
An opportunity like this does not come every day.
\end{quote}

Masculine.

A formal word; everyday speech says \hi{मौक़ा}.

\begin{cousinweb}
\textbf{\hi{अवसर}} (\emph{avsar}) is a tatsama, taken whole from Sanskrit \emph{avasara}, \textquotedblleft{}an occasion, a favourable opportunity\textquotedblright{} (Monier-Williams).
\end{cousinweb}

\begin{grammarlens}[title={What you've built}]
One more word for saying what happens if.
\end{grammarlens}

\subsection*{Your turn}
\begin{itemize}
\raggedright
  \item \emph{Say it:} \emph{avsar}
  \item \emph{Say it:} \emph{avsar}, once more
  \item \emph{Say it:} \emph{Aisā avsar roz nahīṃ ātā.} — An opportunity like this does not come every day.
  \item \emph{Recall:} say the Hindi for a coincidence, then the Hindi for a chance, then say \emph{avsar} again
\end{itemize}

\begin{tcolorbox}[breakable,colback=teal!4,colframe=teal!35!black,title={Before you move on}]
What does \hi{अवसर} mean? (\textquotedblleft{}An opportunity (formal)\textquotedblright{}.) Say it once more.
\end{tcolorbox}
\section[\texorpdfstring{\hi{दाँव} (dā̃v)}{दाँव (dā̃v)}]{\hi{दाँव} (dā̃v) — a turn (in a game); a stake}

\label{lesson:HI-C544-dav}



\begin{quote}
Before the new one: say the Hindi for a chance, then the Hindi for an opportunity (formal).
\end{quote}

\subsection*{You'll want to know: \hi{दाँव}}
\textbf{\hi{दाँव}} — \emph{dā̃v} — \textquotedblleft{}a turn (in a game); a stake\textquotedblright{}.

\begin{quote}
\hi{अगर} \hi{इस} \hi{बार} \hi{मेरा} \hi{दाँव} \hi{लग} \hi{गया}, \hi{तो} \hi{मैं} \hi{जीत} \hi{जाऊँगा।}
\end{quote}

\begin{quote}
\emph{Agar is bār merā dā̃v lag gayā, to maiṁ jīt jāūṅgā.}
\end{quote}

\begin{quote}
If my luck holds this time, I will win.
\end{quote}

Masculine.

\hi{दाँव} \hi{लगना}: one's turn or one's gamble comes off.

\begin{cousinweb}
\textbf{\hi{दाँव}} (\emph{dā̃v}) is a tadbhava, inherited from Sanskrit \emph{dātu}, \textquotedblleft{}a share\textquotedblright{} (CDIAL 6258). CDIAL gives the Hindi word as \textquotedblleft{}a turn\textquotedblright{}: one's share of the play.
\end{cousinweb}

\begin{grammarlens}[title={What you've built}]
One more word for saying what happens if.
\end{grammarlens}

\subsection*{Your turn}
\begin{itemize}
\raggedright
  \item \emph{Say it:} \emph{dā̃v}
  \item \emph{Say it:} \emph{dā̃v}, once more
  \item \emph{Say it:} \emph{Agar is bār merā dā̃v lag gayā, to maiṁ jīt jāūṅgā.} — If my luck holds this time, I will win.
  \item \emph{Recall:} say the Hindi for a chance, then the Hindi for an opportunity (formal), then say \emph{dā̃v} again
\end{itemize}

\begin{tcolorbox}[breakable,colback=teal!4,colframe=teal!35!black,title={Before you move on}]
What does \hi{दाँव} mean? (\textquotedblleft{}A turn (in a game); a stake\textquotedblright{}.) Say it once more.
\end{tcolorbox}
